\section{Hybrid Switching Model and Design Rationale}
\label{sec:switching}

\subsection{The Hybrid Switching Model}

HyNoC's switching model is best described as \emph{circuit-switched wormhole}: it
combines a \textbf{circuit-establishment phase} with a \textbf{wormhole data
phase}.

In the \textbf{circuit-establishment phase}, the ingress port asserts a request to
a specific egress port (identified from the routing flit) and waits for a grant.
This request/grant handshake propagates along the route: once the first router
grants the request, the ingress begins forwarding flits, which in turn trigger the
same handshake at the next router. A dedicated forwarding path through the network
is thus progressively opened, hop by hop.

In the \textbf{wormhole data phase}, payload flits are streamed flit-by-flit along
the established path, regulated by FIFO-level flow control. The path remains
dedicated to this packet until the stop bit is observed, ensuring in-order,
contention-free delivery once the channel is open.

This combination provides the \emph{low buffering} advantage of wormhole switching
(only the ingress FIFO is needed, no full packet storage) while offering the
\emph{path dedication} advantage of circuit switching (no mid-network arbitration
contention during payload transfer). The trade-off is a path-establishment latency
proportional to the route length before the first payload flit is delivered at the
destination.

\subsection{Static Route Management as an Alternative to Virtual Channels}

Virtual channels were introduced primarily to: (1) provide additional logical paths
to avoid deadlock; and (2) reduce head-of-line blocking under congestion. HyNoC
addresses both goals differently.

\textbf{Deadlock avoidance} is achieved through source routing itself: the
complete path is determined by the sender and can be constrained to acyclic
subgraphs (e.g., dimension-order routing) without any router-internal mechanism.

\textbf{Congestion avoidance} is addressed at two levels. First, the compiler or
runtime system assigns routes across multiple available shortest paths, distributing
traffic away from known hotspots. The number of shortest paths between two nodes in
an $N \times M$ mesh is a classical combinatorial result~\cite{knuth1997}: each path consists of a fixed sequence of moves along each dimension, taken in any order, and the count is given by the multinomial coefficient:

\begin{equation}
  p = \binom{(N-1)+(M-1)}{N-1} = \frac{[(N-1)+(M-1)]!}{(N-1)!\,(M-1)!}
  \label{eq:paths2d}
\end{equation}

This generalizes to a $K$-dimensional network $n_1 \times n_2 \times \cdots \times
n_K$ as:

\begin{equation}
  p = \frac{\left[\sum_{i=1}^{K}(n_i-1)\right]!}{\prod_{i=1}^{K}(n_i-1)!}
  \label{eq:pathsnd}
\end{equation}

Note that~(\ref{eq:paths2d}) and~(\ref{eq:pathsnd}) give the number of shortest
paths between the \emph{most distant} node pair (corner-to-corner). For an
arbitrary pair separated by $(d_x, d_y)$ hops the count is $\binom{d_x+d_y}{d_x}$,
which is lower. Table~\ref{tab:paths} reports corner-to-corner values for
representative 2D mesh topologies.

\begin{table}[h]
\centering
\caption{Shortest path counts in 2D mesh NoCs (corner-to-corner)}
\label{tab:paths}
\begin{tabular}{@{}lrrl@{}}
\toprule
Topology & Routers & Shortest paths & Path length \\
\midrule
$2\times 2$   &   4 &               2 &  2 \\
$4\times 4$   &  16 &              20 &  6 \\
$8\times 8$   &  64 &           3\,432 & 14 \\
$16\times 16$ & 256 & 155\,117\,520    & 30 \\
\bottomrule
\end{tabular}
\end{table}

Second, for applications with more dynamic traffic, the network topology can be
enriched with additional relay routers (without local interfaces), increasing path
diversity proportionally to the added area.

\subsection{Area Comparison with Virtual Channels}

Following~\cite{mello2005}, router area scales linearly with the number of VCs
(each VC requires an additional input buffer per port). An $8\times 8$ mesh with
4~VCs consumes the equivalent area of a $16\times 16$ mesh without VCs, yet the
latter offers 66$\times$ more shortest paths between corner nodes. On FPGA, where
block RAMs are scarce, minimizing buffer count is a primary concern.

The VC area model is itself optimistic in two respects: (1) the crossbar logic and
VC scheduler also grow with VC count; (2) VC path diversity is an upper bound,
since in practice the VC is assigned at the source and remains fixed for the entire
path, providing less diversity than a true additional physical dimension.
Conversely, the physical router area model does not capture link wiring cost, which
may be relevant in ASIC contexts.

\subsection{Head-of-Line Blocking and Mitigations}

Once a channel is granted, it remains dedicated until the stop bit of the last
flit is received. A long packet therefore delays other packets competing for the
same egress port. Two complementary mitigations are available in HyNoC:

\begin{enumerate}
  \item \textbf{Network enrichment}: adding more routers distributes traffic
    across more egress ports, reducing the probability that multiple flows compete
    for the same port. The cost is increased end-to-end \emph{latency} due to
    additional hops---this is the primary performance trade-off, not throughput.
  \item \textbf{Multiple local interfaces}: a node can expose several local
    interfaces connected to different router ports, providing as many independent
    logical channels as needed. This replicates the path diversity benefit of VCs
    without duplicating internal router buffers. It is the recommended approach
    when a node requires simultaneous, non-blocking communication with multiple
    peers.
\end{enumerate}

\subsection{Scope and Limitations}

The topology-based approach is most effective when traffic patterns are known at
compile time and routes can be assigned statically. When traffic is entirely
unpredictable at design time and topology enrichment is not feasible, virtual
channels remain the more appropriate solution. HyNoC targets systems---distributed
VLIW processors on FPGA---where compile-time route assignment is natural and FPGA
resource budgets make buffer minimization a primary concern.
